% ============================================================
%  Chapitre 5 — Parallélisation en mémoire partagée (OpenMP)
% ============================================================
\chapter{Parallélisation en mémoire partagée : OpenMP}
\label{chap:openmp}

\section{Stratégie de parallélisation}

\subsection{Identification des boucles parallélisables}

La boucle principale contient deux phases susceptibles d'être parallélisées :

\begin{description}
    \item[Phase fourmis] La boucle sur les $m = 5\,000$ fourmis est \emph{embarrassingly
          parallel} : chaque fourmi avance indépendamment. La seule dépendance partagée
          est la carte de phéromones en lecture (\texttt{m\_map\_of\_pheronome}) et,
          en écriture, le buffer de mise à jour (\texttt{m\_buffer\_pheronome}) via
          \texttt{mark\_pheronome()}. L'écriture concurrente sur des cellules distinctes
          constitue une \emph{race condition tolérable} : l'algorithme ACO est défini
          de façon stochastique et les ordres d'écriture alternatifs produisent des
          résultats équivalents sur un grand nombre d'itérations.
    \item[Phase évaporation] La boucle sur les $(N)^2 = 262\,144$ cellules est
          \emph{totalement indépendante} (pas de dépendance entre cellules).
          Elle constitue le candidat idéal pour la parallélisation.
\end{description}

\subsection{Pragma choisis}

\begin{lstlisting}[language=C++, caption={OpenMP sur la boucle fourmis (ant.cpp)}]
void ant_colony::advance_all(..., std::size_t& cpteur_food) {
    std::size_t local_food = 0;
    #pragma omp parallel for schedule(dynamic, 16) \
                             reduction(+:local_food)
    for (std::size_t i = 0; i < size(); ++i)
        advance_one(i, phen, land, pos_food, pos_nest, local_food);
    cpteur_food += local_food;
}
\end{lstlisting}

Le choix \texttt{schedule(dynamic, 16)} est motivé par la \textbf{variabilité du
nombre de sous-pas} d'une fourmi à l'autre (dépendant du terrain fractal traversé).
Un ordonnancement statique laisserait certains threads inactifs en fin de boucle tandis
que d'autres terminent des fourmis bloquées sur un terrain difficile.
La granule de 16 fourmis par chunk limite le surcoût de gestion du scheduler.

\begin{lstlisting}[language=C++, caption={OpenMP sur l'évaporation (pheronome.hpp)}]
void do_evaporation() {
    #pragma omp parallel for schedule(static)
    for (std::size_t i = 1; i <= m_dim; ++i)
        for (std::size_t j = 1; j <= m_dim; ++j) {
            m_buffer_pheronome[i * m_stride + j][0] *= m_beta;
            m_buffer_pheronome[i * m_stride + j][1] *= m_beta;
        }
}
\end{lstlisting}

Pour l'évaporation, \texttt{schedule(static)} est suffisant car la charge est
parfaitement régulière (même travail par cellule).

\section{Résultats expérimentaux}

\subsection{Mesures de temps par phase}

Les mesures suivantes ont été obtenues en faisant varier le nombre de threads OpenMP
via la variable d'environnement \texttt{OMP\_NUM\_THREADS}, sur la version \texttt{03\_Parallele\_01/}.
Chaque point est la médiane de 3 runs de 15000 itérations, en mode \texttt{--no-render}
pour isoler le coût calculatoire.

\begin{table}[h]
\centering
\caption{Temps par phase en fonction du nombre de threads OpenMP (ms/it)}
\label{tab:omp_timing}
\begin{tabular}{lrrrr}
\hline
\textbf{Threads} & \textbf{Fourmis} & \textbf{Évaporation} & \textbf{Update} & \textbf{Total} \\
\hline
1  & 1.821 & 0.203 & 0.003 & 2.027 \\
2  & 1.029 & 0.103 & 0.003 & 1.137 \\
4  & 0.587 & 0.061 & 0.004 & 0.651 \\
8  & 0.354 & 0.052 & 0.006 & 0.411 \\
12 & 0.259 & 0.051 & 0.006 & 0.315 \\
14 & 0.232 & 0.051 & 0.006 & 0.289 \\
\hline
\end{tabular}
\end{table}

\subsection{Accélération et efficacité}

\begin{table}[h]
\centering
\caption{Accélération $S(p)$ et efficacité $E(p)$ — OpenMP}
\label{tab:omp_speedup}
\begin{tabular}{lrrr}
\hline
\textbf{Threads $p$} & \textbf{$T_p$ (ms/it)} & \textbf{$S(p)$} & \textbf{$E(p)$ (\%)} \\
\hline
1  & 2.027 & 1.00 & 100.0 \\
2  & 1.137 & 1.78 &  89.2 \\
4  & 0.651 & 3.12 &  77.9 \\
8  & 0.411 & 4.93 &  61.6 \\
12 & 0.315 & 6.44 &  53.6 \\
14 & 0.289 & 7.00 &  50.0 \\
\hline
\end{tabular}
\end{table}

% TODO : ajouter figure du graphe speedup (courbe S(p) vs courbe idéale et Amdahl)
% \begin{figure}[h]
%     \centering
%     \includegraphics[width=0.7\textwidth]{Figures/openmp_speedup.png}
%     \caption{Accélération OpenMP en fonction du nombre de threads}
%     \label{fig:omp_speedup}
% \end{figure}

\section{Analyse}

\subsection{Loi d'Amdahl}

La loi d'Amdahl prédit l'accélération maximale pour une fraction parallèle $f$ :

\begin{equation}
    S_{\infty} = \frac{1}{1 - f}
    \label{eq:amdahl}
\end{equation}

À 14 threads, $S(14) = 7.00$. En isolant la partie séquentielle $s = 1 - f$, on obtient :

\[
    s \approx \frac{1}{S(14)} - \frac{1}{14} \approx 0.071
\]

La fraction séquentielle résiduelle est d'environ $7.1\%$ du temps d'exécution.
Elle est principalement constituée de :

\begin{itemize}
    \item la réduction finale du compteur de nourriture (\texttt{reduction(+:local\_food)}) ;
    \item l'opération \texttt{update()} (swap des buffers + mise à jour des conditions limites) ;
    \item le surcoût intrinsèque des \texttt{parallel for} (création/destruction de l'équipe OMP).
\end{itemize}

\subsection{Saturation à partir de 12 threads}

On observe que le gain entre 12 et 14 threads reste modéré ($6.44 \to 7.00$, soit
$\approx +8.8\%$ pour $+17\%$ de threads). Ceci s'explique par la nature hybride de la puce :
les 12 premiers threads occupent les 6 P-cores (2 threads/cœur via HT), tandis que
les 14\textsuperscript{e} et 16\textsuperscript{e} threads mobilisent des E-cores
de cadence inférieure. Par ailleurs, la pression sur le cache L3 partagé (24 Mo)
commence à croître significativement lorsque plusieurs threads accèdent simultanément
à la grille de phéromones de taille $\approx 2 \times 4 \approx 8$ Mo.

\subsection{Évaporation en plateau}

Le temps d'évaporation se stabilise autour de $0.050$--$0.051$~ms/it à partir de
8 threads et ne diminue plus significativement au-delà. Ce plafonnement indique
que la phase d'évaporation est
\textbf{limitée par la bande passante mémoire} (memory-bound) : à 2 threads, le
débit mémoire augmente encore, mais le gain marginal devient rapidement faible pour
ce volume de données ($\approx 4$~Mo de buffer pheromone parcouru séquentiellement).
L'augmentation du nombre de threads au-delà de 8 n'apporte donc qu'un bénéfice limité
sur cette phase spécifique.
